\chapter{Gravity: Holonomy, Torsion, and Collective Response}
\label{chap:gravedad}
\index{gravity}\index{Cartan--Holst}\index{holonomy}
\index{torsion}\index{graviton!non-primitive}

Gravity is not introduced here through a decimal formula for \(G\).
Its construction assembles two already typed branches. In the first, a route preserves
orientation and memory; a connection compares frames at distinct points; its
holonomy measures return and its curvature the closure defect. In the second,
the dimensional chart constructs \(G_{\rm int}(\theta_G)\) and
\(\kappa_{\rm int}\) from the internal action scale, velocity, and
elementary duration. Only after both branches are available is the
Palatini--Cartan--Holst action written. No external metrological figure participates in
that precedence.

\section{Returns: position, orientation, phase, and memory}

Let \(F_{\rm obs}\) be the observable chronology of the transport and
\[
 \mathcal K_{27}=F_{\rm obs}^{27}.
\]
The exact hierarchy distinguishes three returns:
\[
 \begin{array}{c|c}
 27\ \text{steps}&
 \text{position return and orientation reversal},\\
 54\ \text{steps}&
 \text{position and orientation return},\\
 108\ \text{steps}&
 \text{complete return of the observable phase}.
 \end{array}
\]
The observable return does not erase the accumulated record. If \(g(e)\) counts
an orientation reversal and \(s(e)\) an effective sheet change, then
\[
 c(\gamma)=\sum_{e\subset\gamma}(g(e),s(e))
\]
satisfies the cocyclic law of concatenation. In the canonical blocks,
\[
 c_{27}=(1,18),\qquad c_{108}=(4,72).
\]
Therefore, the observable projection of \(\mathcal K_{27}\) has order
four, whereas the cumulative coordinate continues to grow.  Returning to the
same phase does not mean returning to the same history.

\begin{theorem}[Cumulative coverage of return]
\label{thm:cobertura-retorno-grav-rev6}
Over \(\mathcal O_{108}\times\mathbb N^2\), the lift
\[
 \widetilde F(x,m)=
 \bigl(F_{\rm obs}x,m+c(x\to F_{\rm obs}x)\bigr)
\]
satisfies
\[
 \widetilde{\mathcal K}_{27}(x,m)
 =\bigl(\mathcal K_{27}x,m+(1,18)\bigr)
\]
and
\[
 \widetilde{\mathcal K}_{27}^{\,4}(x,m)
 =\bigl(x,m+(4,72)\bigr)\neq(x,m).
\]
\end{theorem}

\begin{proof}
The constancy of \(c_{27}\) along the orbit and the concatenation law give
\[
 c_{108}=\sum_{j=0}^{3}c_{27}(\mathcal K_{27}^jx)
 =4(1,18).
\]
The first coordinate returns; the second preserves the accumulated memory.
\end{proof}

This is the datum that to geometric realization must preserve.  One does not
identify the discrete memory with torsion: one constructs to representation
that allows them to is compared.

\section{From the path groupoid to the Cartan group}

Let \(\mathcal G_{\rm mem}\) be the groupoid of enriched routes, with
composition by concatenation and inverse by reversal. Fix an element
\(R_4\in\operatorname{Spin}^+(1,3)\) of order four and a unit temporal vector
\(u\in\mathbb R^{1,3}\) fixed by \(R_4\). If
\[
c(\gamma)=\bigl(c_g(\gamma),c_s(\gamma)\bigr)\in\mathbb Z^2
\]
is the additive cocycle of inversion and sheet change, we define
\[
\rho_{\rm C}^{\rm disc}(\gamma)
=
\left(
R_4^{\,c_g(\gamma)},
\ell_0c_s(\gamma)u
\right)
\in\operatorname{Spin}^+(1,3)\ltimes\mathbb R^{1,3}.
\]

\begin{theorem}[Discrete realization of the return cocycle]
\label{thm:realizacion-discreta-cartan-rev6}
\(\rho_{\rm C}^{\rm disc}\) is a functor from the route groupoid to the
Cartan group. In particular, it preserves identity, concatenation, and inversion, and
transports the return \(27\to54\to108\) without erasing the cumulative coordinate.
\end{theorem}

\begin{proof}
The semidirect product law is
\[
(R,a)(S,b)=(RS,a+Rb).
\]
Since \(R_4u=u\) and \(c\) is additive,
\[
\begin{aligned}
\rho_{\rm C}^{\rm disc}(\gamma_2)
\rho_{\rm C}^{\rm disc}(\gamma_1)
&=
\left(
R_4^{c_g(\gamma_2)+c_g(\gamma_1)},
\ell_0\bigl(c_s(\gamma_2)+c_s(\gamma_1)\bigr)u
\right)\\
&=\rho_{\rm C}^{\rm disc}(\gamma_2\gamma_1).
\end{aligned}
\]
The identity route has zero cocycle, and reversal changes its sign, so
identity and inverse are preserved. The return formulas follow from
\(c_{27}=(1,18)\) and \(c_{108}=4c_{27}\).
\end{proof}

This theorem constructs the discrete bridge. A subsequent smooth realization consists of a representation
\[
 \rho_{\rm C}:\mathcal G_{\rm mem}
 \longrightarrow\operatorname{Spin}(1,3)\ltimes\mathbb R^{1,3}
\]
and in fields \(e^I,\omega^{IJ}\) whose holonomy intertwines that representation.
The diagram fixing the type is
\[
 \operatorname{Hol}_{\mathcal A_{\ell_0}}
   \bigl(\mathfrak R_{\rm C}(\gamma)\bigr)
 =
 \rho_{\rm C}\bigl(\operatorname{Hol}_{\rm TPK}(\gamma)\bigr).
\]
Identities, inversion, concatenation, and subdivision must be preserved.  This
equation does not rename the carry as curvature: it requires the already represented
discrete transport and the geometric transport to realize the same composition.

Once to positive scale is fixed \(\ell_0\), the joint connection is
\[
 \mathcal A_{\ell_0}
 =
 \begin{pmatrix}
 \omega&\ell_0^{-1}e\\
 0&0
 \end{pmatrix}.
\]

\begin{proposition}[Cartan Joint Curvature]
\label{prop:curvatura-conjunta-cartan-rev6}
The curvature of \(\mathcal A_{\ell_0}\) is
\[
 \mathcal F_{\ell_0}
 =\dd\mathcal A_{\ell_0}
  +\mathcal A_{\ell_0}\wedge\mathcal A_{\ell_0}
 =
 \begin{pmatrix}
 F_\omega&\ell_0^{-1}T_{\rm tor}\\
 0&0
 \end{pmatrix},
\]
where
\[
 F_\omega^{IJ}
 =\dd\omega^{IJ}+\omega^I{}_K\wedge\omega^{KJ},
 \qquad
 T_{\rm tor}^{I}=D_\omega e^I.
\]
\end{proposition}

\begin{proof}
The matrices of forms are multiplied using the exterior product.  The
diagonal block produces \(F_\omega\); the translation block produces
\(\dd e+\omega\wedge e=D_\omega e\).
\end{proof}

Curvature and torsion are two components of the same transport defect,
but they perform distinct functions.  Curvature measures frame return;
torsion measures closure of the co-tetrad.

\section{Gauss--Stokes and areal capacity}

In a finite cellular complex, for every \(1\)-cochain \(\Omega\) and every
\(2\)-chain \(z\),
\[
 \boxed{\quad
 \langle\delta\Omega,z\rangle
 =\langle\Omega,\partial z\rangle.
 \quad}
\]
The identity is the definition of the coboundary as the adjoint of the boundary.
When \(z\) sums the coherently oriented faces of a
pseudomanifold, the interior edges cancel, and only the geometric
boundary remains.  The same equality has already been read as
vorticity--circulation and as curvature--holonomy.

The areal scale also has an exact model.  In the cubic graph
\(N\times N\times N\), every cut between two opposite faces contains at least
\(N^2\) edges, and to transverse layer attains the bound.  For \(N=9\),
\[
 9^3=729=27^2.
\]
This cardinal equality relates volume and minimal cut; it does not identify the
groups \((\mathbb Z/9)^3\) and \((\mathbb Z/27)^2\), which have distinct
exponents.

\section{The Palatini--Cartan--Holst Action}

Let \(M\) be a smooth oriented manifold of dimension four, with nondegenerate
co-tetrad \(e^I\), Lorentz connection
\(\omega^{IJ}=-\omega^{JI}\), and boundary conditions that make the boundary
variational term vanish.  Let us write
\[
 \Sigma^{IJ}=e^I\wedge e^J,\qquad
 (\star X)^{IJ}=\frac12\epsilon^{IJ}{}_{KL}X^{KL},
 \qquad
 G_{\rm int}=G_{\rm int}(\theta_G),\qquad
 \kappa_{\rm int}=\frac{8\pi G_{\rm int}}{c_{\rm int}^{4}}.
\]
These last two quantities are not defined through the action: they follow from
the constitutive construction of \cref{eq:G-interno}. The
\cref{thm:homogeneidad-pch} first proves that the following combination
belongs to the action line.
The first-order action is
\[
\boxed{
 S[e,\omega,\Psi]
 =\frac1{2\kappa_{\rm int}c_{\rm int}}\int_M
 \Sigma_{IJ}\wedge
 \left(\star F^{IJ}+\frac1\gamma F^{IJ}\right)
 -\frac{\Lambda_{\rm int}}{\kappa_{\rm int}c_{\rm int}}
  \int_M\operatorname{vol}_e
 +S_{\rm m}[e,\omega,\Psi].}
\]
The internal parameter constructed in the previous chapter specializes
\[
 \gamma=\Gamma_{6\to8}.
\]
The substitution is performed after obtaining
\(\Gamma_{6\to8}\) from \(A,C^\ast,q_\pm\) and the hexad--octad inclusion; the action is not used to redefine those inputs.

\begin{theorem}[Cartan--Holst Equation]
\label{thm:cartan-holst-rev6}
If the spin current is normalized through
\[
 \delta_\omega S_{\rm m}
 =-\frac12\int_M\delta\omega^{IJ}\wedge\sigma_{IJ},
\]
the variation with respect to \(\omega\) produces
\[
 \boxed{\quad
 D_\omega
 \left[(\star+\gamma^{-1}I)\Sigma^{IJ}\right]
 =\kappa_{\rm int}c_{\rm int}\,\sigma^{IJ}.
 \quad}
\]
\end{theorem}

\begin{proof}
One uses \(\delta F^{IJ}=D_\omega\delta\omega^{IJ}\), integrates by parts, and
sets the coefficient of the arbitrary variation equal to zero.  The cosmological
term does not depend on \(\omega\).
\end{proof}

In Lorentzian signature, \(\star^2=-I\) on bivectors.  For
\(\gamma\in\mathbb R\setminus\{0\}\),
\[
 (\star+\gamma^{-1}I)^{-1}
 =\frac{\gamma^2}{1+\gamma^2}
  (-\star+\gamma^{-1}I).
\]
Therefore, if the spin current vanishes,
\[
 T_{\rm tor}^I=0,\qquad
 \omega=\mathring\omega(e).
\]
With spinorial matter, torsion is determined algebraically by the
current.  In the minimal action, it does not constitute an independent
propagating degree of freedom.

\section{Metric equation, spin, and geometric balance}

To fix the normalization unambiguously, we defines the sources measured on
the action line by
\[
 \delta_g S_{\rm m}
 =-\frac12\int_M\operatorname{vol}_e\,
 \mathcal T^{S}_{\mu\nu}\,\delta g^{\mu\nu}.
\]
When the tetrad coordinates have the dimension of length, we write
\[
 T^{\rm phys}_{\mu\nu}:=c_{\rm int}\mathcal T^{S}_{\mu\nu},
 \qquad
 U^{\rm phys,spin}_{\mu\nu}
 :=c_{\rm int}\mathcal U^{S,\rm spin}_{\mu\nu}.
\]
After algebraically eliminating the connection, the metric equation adopts the two equivalent forms
\[
 \boxed{
 G_{\mu\nu}+\Lambda_{\rm int}g_{\mu\nu}
 =\kappa_{\rm int}c_{\rm int}
 \bigl(\mathcal T_{\mu\nu}^{S,\rm LC}
       +\mathcal U_{\mu\nu}^{S,\rm spin}\bigr)
 =\kappa_{\rm int}
 \bigl(T_{\mu\nu}^{\rm phys,LC}
       +U_{\mu\nu}^{\rm phys,spin}\bigr).}
\]
The spin contribution is quadratic in the corresponding current and
depends on the matter coupling; its coefficient is not fixed without declaring
the fermionic action and the duality conventions.

The geometric identities
\[
 D_\omega T_{\rm tor}^{I}=F^{I}{}_{J}\wedge e^{J},
 \qquad
 D_\omega F^{IJ}=0
\]
imply, after removing the torsion,
\[
 \nabla^\mu\bigl(T_{\mu\nu}^{\rm phys,LC}
             +U_{\mu\nu}^{\rm phys,spin}\bigr)=0.
\]
This Bianchi--Noether balance is to consistency condition for the
gravitational transducer: not additional term can is incorporated without
proceeding from to compatible action or satisfying the same balance.

\section{Area--volumetric correspondence \(\pi/\varphi^2\)}

The calendar \((\Sigma,\Pi,\Pi)\) produces the incidence matrix
\[
 F_{\rm av}=\begin{pmatrix}0&1\\1&1\end{pmatrix}.
\]
Its Parry measure gives
\(\mu_{\rm ar}/\mu_{\rm vol}=\varphi^{-2}\), and the regional reader marks the
closure exactly once \(\pi_{\rm HMT}\).  The marked return satisfies
\[
 \Theta_n
 =\pi_{\rm HMT}\frac{F_{n-1}}{F_{n+1}}
 \longrightarrow
 \frac{\pi_{\rm HMT}}{\varphi^2}.
\]
This ratio is not another version of \(G\). It is the dimensionless correspondence between
the areal and volumetric readings. Under the same positive amplitude
\(\mathcal M\),
\[
 \mathcal I_{\rm av}=\mu_{\rm vol}\mathcal M,\qquad
 \mathcal G_{\rm av}=\pi_{\rm HMT}\mu_{\rm ar}\mathcal M
\]
give exactly
\[
 \frac{\mathcal G_{\rm av}}{\mathcal I_{\rm av}}
 =\frac{\pi_{\rm HMT}}{\varphi^2}.
\]
On the flat sheet, the diagonal restriction is
\(\mathcal G_{\rm tor}=\mathcal I_{\rm tor}\).  Local equivalence and
boundary ambivalence are therefore two regimes of the reader, not two
incompatible values of a mass.

\section{Typing of \(G\) as gravitational transducer and causal order of its readers}

The type and homogeneity of \(G\) belong to the prior constitutive
construction. The \cref{eq:G-interno} fixes the covariant family
\(G_{\rm int}(\theta_G)\) once and for all, and the
\cref{thm:homogeneidad-pch} proves that its insertion into the action has
the dimension of action. This chapter does not redefine that family: it uses its
section together with the holonomy already constructed. Causality is not linear,
but convergent:
\[
\begin{aligned}
 \mathcal B_{\rm geom}&:
 \text{return and cocycle}\to\text{Cartan representation}
 \to(e,\omega),\\
 \mathcal B_{\rm dim}&:
 (\hbar_{\rm int},c_{\rm int},t_0,\theta_G)
 \to G_{\rm int}\to\kappa_{\rm int},
\end{aligned}
\]
\[
 (\mathcal B_{\rm geom},\mathcal B_{\rm dim})
 \to S_{\rm PCH}
 \to\text{algebraic torsion, metric equation and balance}.
\]
The metrological comparison belongs to the end and does not select any of
the two branches.

\Needspace{22\baselineskip}
\section{Collective origin of the gravitational mode in the holonomy of paths}

The basic objects of the preceding action are the co-tetrad and the connection.
Cartan torsion is algebraic; curvature has the propagating sector that
contains gravitational waves. Accordingly, the formalism does not need to
postulate an elementary graviton in order to define the interaction. A possible
quantization of the perturbations may produce a collective mode of
helicities \(\pm2\); that mode would be an excitation of geometric
transport, not a datum added at the level of APP, TRIT, or TPK.

This distinction also excludes a frequent confusion. The existence of
classical waves proves propagation of curvature; by itself it does not prove
a fundamental quantum. Conversely, describing the field geometrically does not
eliminate its radiative degrees of freedom. HMT--MD locates the question properly:
memory and its holonomy are constructed first; the spectrum of their excitations
is then studied.
